\documentclass[11pt,a4paper]{article}
\usepackage[margin=1in]{geometry}
\usepackage{amsmath,amssymb,amsfonts}
\usepackage{booktabs}
\usepackage{array}
\usepackage{hyperref}
\usepackage{xcolor}

\hypersetup{colorlinks=true,linkcolor=blue,citecolor=blue,urlcolor=blue}

\newcommand{\CP}{\mathbb{CP}}

\begin{document}

\title{Reconciliation of Two DFD Fermion Mass Scripts:\\
\texttt{compute\_all\_masses.py} vs.\ \texttt{compute\_mass\_exponents.py}}

\author{Research Note for DFD v108}

\date{March 2026}

\maketitle

%==========================================================================
\section{Summary}
%==========================================================================

Two Python scripts exist in the DFD repository that compute fermion masses
from the formula $m_f = A_f\,\alpha^{n_f}\,(v/\sqrt{2})$.  They use
\textbf{different prefactors and different exponents} for several fermions.
This note identifies which is canonical and explains the origin of the
discrepancy.

\textbf{Conclusion:} They are \textbf{genuinely different models}, not
different conventions for the same model.  Script~A matches the verified
v67 Mass Dictionary and is canonical.  Script~B is an exploratory variant
that should not be used for publication.

%==========================================================================
\section{Side-by-Side Comparison}
%==========================================================================

\begin{table}[h]
\centering
\renewcommand{\arraystretch}{1.15}
\begin{tabular}{@{}l cc cc@{}}
\toprule
& \multicolumn{2}{c}{\textbf{Script A} (\texttt{compute\_all\_masses.py})}
& \multicolumn{2}{c}{\textbf{Script B} (\texttt{compute\_mass\_exponents.py})} \\
\cmidrule(lr){2-3} \cmidrule(lr){4-5}
Fermion & $A_f$ & $n_f$ & $A_f$ & $n_f$ \\
\midrule
$e$   & $2/3$   & 2.5 & $2/\pi$      & 2.5 \\
$\mu$ & $1$     & 1.5 & $1$          & 1.5 \\
$\tau$& $\sqrt{2}$ & 1.0 & $\sqrt{2}$ & 1.0 \\
$u$   & $8/3$   & 2.5 & $2\sqrt{2}$  & 2.5 \\
$c$   & $1$     & 1.0 & $1$          & 1.0 \\
$t$   & $1$     & 0   & $1$          & 0   \\
$d$   & $6$     & \textbf{2.5} & $1/2$  & \textbf{2.0} \\
$s$   & $6/7$   & 1.5 & $\sqrt{3}/2$ & 1.5 \\
$b$   & $1/42$  & \textbf{0}   & $\pi$  & \textbf{1.0} \\
\midrule
Mean $|$error$|$ & \multicolumn{2}{c}{1.42\%} & \multicolumn{2}{c}{1.73\%} \\
\bottomrule
\end{tabular}
\caption{Comparison of the two scripts.  Bold entries highlight where
exponents differ.  Five of nine prefactors also differ.}
\label{tab:comparison}
\end{table}

%==========================================================================
\section{Are These the Same Model?}
\label{sec:same}
%==========================================================================

\textbf{No.}  They differ in two structurally significant ways:

\subsection{Different Exponents for $b$ and $d$}

\begin{itemize}
\item \textbf{Bottom quark:}
  Script~A has $n_b = 0$ (the $b$ quark sits at the Higgs vertex, same as $t$,
  with $k_f = 1$, $k_H = +1$, and QCD dressing shifting $n = 1 \to 0$).
  Script~B has $n_b = 1$ (using a ``linear rule'' $k_f = 4 - g$ with $g = 3$
  giving $k_f = 1$, $k_H = +1$, $n = 1$, and \emph{no} QCD dressing shift).

\item \textbf{Down quark:}
  Script~A has $n_d = 5/2$ (same as the electron and up quark---all
  first-generation fermions at maximum CP$^2$ distance).
  Script~B has $n_d = 2$ (using $k_f = 3$, $k_H = +1$, giving $(3+1)/2 = 2$).
\end{itemize}

These differences mean the two scripts assign different bundle degrees $k_f$
to the $b$ and $d$ quarks.  This is a physical disagreement about which
spin$^c$ zero mode each fermion occupies on $\CP^2$.

\subsection{Different Prefactors for 5 Fermions}

\begin{center}
\begin{tabular}{lll}
\toprule
Fermion & Script A & Script B \\
\midrule
$e$   & $2/3$ (rational)      & $2/\pi$ (transcendental) \\
$u$   & $8/3$ (rational)      & $2\sqrt{2}$ (algebraic) \\
$d$   & $6$ (integer)         & $1/2$ (rational) \\
$s$   & $6/7$ (rational)      & $\sqrt{3}/2$ (algebraic) \\
$b$   & $1/42$ (rational)     & $\pi$ (transcendental) \\
\bottomrule
\end{tabular}
\end{center}

Script~A's prefactors are all rational numbers (or $\sqrt{2}$ for $\tau$),
and each can be traced to an explicit operator product: $A_f = G[g] \times$
(sector factor).

Script~B's prefactors include $\pi$ and $2/\pi$, which cannot arise from
the operator algebra of finite groups and matrices.  These are fitting
coefficients, not derived quantities.

%==========================================================================
\section{Which Is Correct?}
\label{sec:correct}
%==========================================================================

\textbf{Script~A} is correct.  The evidence:

\begin{enumerate}
\item \textbf{Matches the v67 dictionary.}
  The v67 Mass Dictionary was produced in January 2026, verified by running
  \texttt{dfd\_v67\_mass\_dictionary\_verify.py}, and explicitly approved.
  Script~A reproduces it exactly.

\item \textbf{Lower mean error.}
  Script~A achieves 1.42\% vs.\ Script~B's 1.73\%.

\item \textbf{Prefactors are derivable.}
  Script~A's prefactors all come from the explicit Yukawa operator:
  \[
  A_f = G[g_f] \times D_\ell[g_f] \quad\text{(leptons)}, \qquad
  A_f = G[g_f] \times Q_d[g_f] \times R_d[g_f] \quad\text{(down quarks)},
  \]
  etc.  Every factor is a matrix element of an algebraically defined operator.
  Script~B's $\pi$ and $2/\pi$ have no such derivation.

\item \textbf{Consistent exponent pattern.}
  In Script~A, all first-generation fermions ($e$, $u$, $d$) share $n_f = 5/2$
  (maximum CP$^2$ distance).  All third-generation fermions ($\tau$, $t$, $b$)
  have $n_f \in \{0, 1\}$ (at or near the Higgs vertex).  This is a
  cleaner, more symmetric pattern.
\end{enumerate}

%==========================================================================
\section{Origin of Script~B}
\label{sec:origin}
%==========================================================================

Script~B (\texttt{compute\_mass\_exponents.py}) was written as an
\emph{exploratory} exercise to derive the exponents $n_f$ from three
``generation rules'' (doubling, linear, triangular) applied to the
bundle degrees.  The exponent derivation in Script~B is valuable---it
provides a candidate mechanism for \emph{why} specific $k_f$ values
arise---but it produced a \textbf{different} set of $(k_f, A_f)$ for
the $b$ and $d$ quarks.

The key discrepancy is that Script~B's ``linear rule'' for down quarks
gives $k_d = 3$ (first generation) and $k_b = 1$ (third generation)
\emph{without} the QCD dressing shift that, in Script~A/v67, moves
$n_b$ from 1 to 0.  Script~B then compensates by using entirely different
prefactors ($A_d = 1/2$ instead of $6$; $A_b = \pi$ instead of $1/42$)
to fit the observed masses.

\subsection{What Script~B Got Right}

\begin{itemize}
\item The lepton and up-quark exponents match Script~A exactly.
\item The exponent derivation rules (doubling for leptons, triangular for
  up quarks) are a valid proposal for the underlying mechanism.
\item The $\mu$, $\tau$, $c$, $t$ entries are identical in both scripts.
\end{itemize}

\subsection{What Script~B Got Wrong}

\begin{itemize}
\item The down-quark sector exponents ($n_d$, $n_b$) disagree with the
  verified v67 dictionary.
\item The prefactors $A_b = \pi$ and $A_e = 2/\pi$ are transcendental
  numbers that cannot come from the finite Yukawa operator algebra.
\item The higher mean error (1.73\% vs.\ 1.42\%) confirms the mismatch.
\end{itemize}

%==========================================================================
\section{Recommendation}
\label{sec:recommendation}
%==========================================================================

\begin{enumerate}
\item \textbf{Use Script~A values (= v67 dictionary) for all publications.}
  The canonical script is now \texttt{mass\_canonical.py}.

\item \textbf{Retire Script~B as an exploratory tool.}
  The exponent derivation mechanism in Script~B (generation rules) is
  interesting but needs revision in the down-quark sector to match the
  canonical values.

\item \textbf{Do not mix prefactors/exponents between the two scripts.}
  They are coupled: changing an exponent requires changing the corresponding
  prefactor to maintain agreement with experiment.

\item \textbf{The free parameter count is ONE, not zero}, when using
  $v/\sqrt{2} = 174.1$~GeV from $G_F$.  The zero-parameter claim requires
  the derived formula $v = M_P\alpha^8\sqrt{2\pi}$, which is at Tier~1.5
  rigor (not yet published-theorem grade).
\end{enumerate}

%==========================================================================
\section{Canonical Values (for Reference)}
\label{sec:canonical}
%==========================================================================

The following table is the single authoritative mass dictionary.

\begin{table}[h]
\centering
\renewcommand{\arraystretch}{1.2}
\begin{tabular}{@{}lccrrc@{}}
\toprule
Fermion & $A_f$ & $n_f$ & $m_{\text{pred}}$ (MeV) & $m_{\text{PDG}}$ (MeV) & Error \\
\midrule
$e$      & $2/3$       & 2.5 &    0.528 &    0.511 & $+3.32\%$ \\
$\mu$    & $1$         & 1.5 &  108.53  &  105.66  & $+2.72\%$ \\
$\tau$   & $\sqrt{2}$  & 1.0 & 1796.7   & 1776.86  & $+1.12\%$ \\
$u$      & $8/3$       & 2.5 &    2.112 &    2.16  & $-2.23\%$ \\
$c$      & $1$         & 1.0 & 1270.5   & 1270     & $+0.04\%$ \\
$t$      & $1$         & 0   & 174100   & 172760   & $+0.78\%$ \\
$d$      & $6$         & 2.5 &    4.752 &    4.67  & $+1.75\%$ \\
$s$      & $6/7$       & 1.5 &   93.03  &   93.0   & $+0.03\%$ \\
$b$      & $1/42$      & 0   & 4145.2   & 4180     & $-0.83\%$ \\
\midrule
\multicolumn{5}{@{}l}{Mean $|$error$|$} & $1.42\%$ \\
\bottomrule
\end{tabular}
\caption{Canonical v67 mass dictionary.  $\alpha = 1/137.036$,
$v/\sqrt{2} = 174.1$~GeV.  One free parameter (the global scale).}
\end{table}

\end{document}
